% =============================================================================
% §5 Conclusion (~0.3 page)
% =============================================================================
\section{Conclusion}
\label{sec:conclusion}

SkillCorpus curates $\sim$$821{,}000$ crawled community skill files
into a corpus of $96{,}401$ skills, scored along utility,
robustness, and safety and matched to tasks by a fine-tuned
retrieval-and-selection stack.  Across our cross-benchmark,
cross-harness, multi-backbone evaluation
(\S\ref{sec:experiments}), integrating the corpus yields
consistent positive gains on all three benchmarks.  They
are large on SkillsBench and modest on GDPVal and QwenClawBench,
strongest on the Raven harness, with per-domain effects that
remain heterogeneous.  Where coverage is thin, growing the corpus
by skill generation is the natural next step.  The 16-class
taxonomy steers benchmark-aware retrieval.  The composite quality score
serves curation and ranking rather than per-task prediction, and
its safety facet is the one component that both gates unsafe
skills and carries a suggestive per-task signal.  
% We intend SkillCorpus to serve as both an open resource and a reusable framework for community skill ecosystems.
We release SkillCorpus as an open, reusable framework for community skill ecosystems; to our knowledge, it is the first end-to-end account of when a curated community corpus improves real agent tasks, and where it does not.
\ifarxiv\else
Our evaluation prioritises breadth over depth, so on the
high-baseline LLM-judged benchmarks the positive effect is
established at the pooled level rather than cell by cell.  The
English-dominant corpus, the two-harness scope, and the 2026-Q2
snapshot bound external validity, leaving broader languages,
harnesses, and human-validated quality labels to future work.
\fi
